% ============================================================
% INTRODUÇÃO (Cap. 1) — a escrever ao final. NÃO incluída no PDF.
% Contém apenas a introdução propriamente dita. Contexto, Objetivos,
% Perguntas de pesquisa e Contribuições foram movidos para um capítulo
% próprio, compilado, em corpo.tex.
% ============================================================

\chapter{Introdução}
\label{cap:intro}
% =====================================================================

As Redes Sociais Digitais (RSD) tornaram-se um substrato central de observação
do comportamento social, e sobre elas se pratica um amplo repertório de análises
computacionais --- da detecção de comunidades à análise de sentimento. Toda essa
prática, porém, repousa sobre uma etapa anterior e raramente problematizada: a
coleta dos dados ---  e, em particular, o seu volume. Quanto se
coleta é quase sempre decidido por conveniência (o que a API entrega, uma janela
de tempo, os itens ``top''), e quase nunca se pergunta se coletar mais
mudaria a conclusão, ou se o volume coletado já bastava. Essa pergunta ---
coletar pouco ou muito muda a conclusão? --- não é exclusiva das RSD; é
uma instância concreta de uma preocupação metodológica mais ampla, presente
sempre que se analisa um recorte de um universo grande demais para ser
observado por inteiro. É por isso que esta dissertação toma as RSD como
caso de estudo: um domínio em que a questão do volume é ao mesmo tempo
aguda, recorrente e, sobretudo, mensurável.

Em trabalho anterior, foi proposta uma especificação conceitual para RSD, com
esquemas conceituais específicos e um esquema genérico que viabiliza armazenar e
consultar dados de plataformas distintas de maneira
consistente~\cite{salgueiro2023,salgueiro2022dbmodels}. Aquele trabalho assume,
como ponto de partida, os dados já coletados. Esta dissertação
problematiza a etapa anterior --- a coleta, e sobretudo o seu volume ---
e o seu efeito sobre as conclusões.

Toda análise computacional sobre uma RSD é executada não sobre o universo de
interações da plataforma, mas sobre um recorte: uma amostra, uma janela
temporal, um conjunto de palavras-chave, os itens mais votados que uma API entrega,
ou um conjunto de dados pré-existente. A questão é se esse recorte, e o volume que
ele descarta, é inocente. Para grafos, sabe-se que amostrar uma rede pode
distorcer fortemente suas propriedades estruturais~\cite{leskovec2006sampling};
resta saber quão sensível é cada tipo de análise efetivamente praticado na
literatura de RSD, e o que acontece quando se coleta mais do que os
trabalhos coletaram. A resposta, obtida no caso das RSD por serem um domínio em
que o efeito é mensurável, informa diretamente a preocupação metodológica de que
ele é instância.

\section{Organização do texto}

O Capítulo~\ref{cap:experimento} apresenta a metodologia --- a lógica,
os eixos de expansão, as métricas, o conjunto de casos fixado e o protocolo
operacional. O Capítulo~\ref{cap:conclusao} resume o estado atual e os próximos
passos. O levantamento que fundamenta a motivação está detalhado no artigo
companion~\cite{barreto2026survey}.

